\documentclass[14pt,a4paper]{extarticle}
\usepackage[T2A]{fontenc}
\usepackage[utf8]{inputenc}
\usepackage[english,russian]{babel}
\usepackage{lmodern}
\usepackage{geometry}
\usepackage{array}
\usepackage{tabularx}
\usepackage{longtable}
\usepackage{ltablex}
\usepackage{booktabs}
\usepackage{hyperref}
\usepackage{indentfirst}
\usepackage{microtype}
\usepackage{graphicx}
\usepackage{setspace}
\usepackage{listings}
\usepackage{graphicx}
\usepackage{caption}
\usepackage{float}
\usepackage{tocloft}
\usepackage{titlesec}
\usepackage{enumitem}
\usepackage{newtxtext,newtxmath}


\titleformat{\section}        % форматируем \section
  {\normalfont\bfseries}      % стиль: обычный + жирный
  {\thesection.}              % номер секции с точкой
  {2mm}                       % расстояние между номером и текстом
  {}                           % перед заголовком можно добавить что-то
\titlespacing*{\section}{1.5cm}{*3}{*1}  % {отступ слева}{отступ сверху}{отступ снизу}

\titleformat{\subsection}        % форматируем \section
  {\normalfont\bfseries}      % стиль: обычный + жирный
  {\thesubsection.}              % номер секции с точкой
  {2mm}                       % расстояние между номером и текстом
  {}                           % перед заголовком можно добавить что-то

\titleformat{\subsubsection}        % форматируем \section
  {\normalfont\bfseries}      % стиль: обычный + жирный
  {\thesubsubsection.}              % номер секции с точкой
  {2mm}                       % расстояние между номером и текстом
  {}                           % перед заголовком можно добавить что-то

\titleformat{\item}        % форматируем \section
  {\normalfont\bfseries}      % стиль: обычный + жирный
  {\theitem.}              % номер секции с точкой
  {2mm}                       % расстояние между номером и текстом
  {}                           % перед заголовком можно добавить что-то
\titlespacing*{\section}{1.5cm}{*3}{*1}  % {отступ слева}{отступ сверху}{отступ снизу}
\titlespacing*{\subsection}{1.5cm}{*3}{*1}  % {отступ слева}{отступ сверху}{отступ снизу}
\titlespacing*{\subsubsection}{1.5cm}{*3}{*1}  % {отступ слева}{отступ сверху}{отступ снизу}
\setlist[enumerate]{
    leftmargin=2.5cm,
    itemindent=0cm
}
\captionsetup{justification=centering}
\renewcommand{\ttdefault}{cmtt}
\keepXColumns
\setlength{\parindent}{3em}
\newgeometry{left=2.5cm,right=2.5cm,top=3cm,bottom=2.5cm}

\renewcommand{\contentsname}{\hspace*{4em}{\smallСодержание}}
\setlength{\cftbeforesecskip}{0.5em}
\renewcommand{\cftsecleader}{\cftdotfill{\cftdotsep}}
\renewcommand{\cftsecaftersnum}{.\hspace{2mm}}

\begin{document}

\setstretch{1.5}

% ----------------------------------------------------------------------
% Титульный лист
% ----------------------------------------------------------------------
\noindent
\begin{minipage}[t]{0.01\textwidth}
    \includegraphics{~/programmer/++Programmer/5_1/LSTU.jpg}
\end{minipage}%
\begin{minipage}[c]{1.1\linewidth}
    \centering
    {\small
    МИНИСТЕРСТВО НАУКИ И ВЫСШЕГО ОБРАЗОВАНИЯ РФ \\[1em]
    ФЕДЕРАЛЬНОЕ ГОСУДАРСТВЕННОЕ БЮДЖЕТНОЕ \\
    ОБРАЗОВАТЕЛЬНОЕ УЧРЕЖДЕНИЕ ВЫСШЕГО ОБРАЗОВАНИЯ \\
    «ЛИПЕЦКИЙ ГОСУДАРСТВЕННЫЙ ТЕХНИЧЕСКИЙ УНИВЕРСИТЕТ»
    }
\end{minipage}%

\vspace{3em}

\noindent
\hspace{1em}%
\begin{minipage}[t]{0.8\textwidth}
\textbf{Институт} \hspace{2em} компьютерных наук \\
\textbf{Кафедра} \hspace{2em} автоматизированных систем управления
\end{minipage}

\vspace{6em}

\begin{center}
    ОТЧЕТ ПО ПЕДАГОГИЧЕСКОЙ ПРАКТИКЕ \\ В ЛГТУ
\end{center}

\vspace{6em}

\noindent
\begin{tabular*}{\textwidth}{@{\extracolsep{\fill}} l c r }
    Студент М-РИТ-25-1 & \makebox[4cm]{\hrulefill} & Станиславчук С.М. \\[2em]
Руководитель практики& \\
старший преподаватель & \makebox[4cm]{\hrulefill} & Качановский Ю.П. \\[2em]
Научный руководитель & \\
доцент& \makebox[4cm]{\hrulefill} & Гаев Л.В. \\
\end{tabular*}

\vfill
\begin{center}
    Липецк, 2026
\end{center}

\thispagestyle{empty}

\newpage
\tableofcontents
\newpage

% ======================================================================
\section*{Введение}
\addcontentsline{toc}{section}{Введение}

Учебная (педагогическая) практика является важнейшим этапом в подготовке магистранта, направленным на закрепление теоретических знаний и формирование практических навыков в области педагогики. Это обязательный компонент основной образовательной программы, нацеленный на развитие профессиональных компетенций будущего преподавателя, овладение основами педагогического мастерства, а также умениями самостоятельного ведения учебно-воспитательной работы.

Основная идея практики заключается в непосредственном погружении в педагогическую деятельность, формировании умений, связанных с проектированием, конструированием и организацией учебного процесса для студенческой аудитории. В ходе практики магистрант развивает стратегическое мышление, коммуникативные навыки и умение эффективно взаимодействовать с учебной группой.

Целью учебной практики является формирование и развитие первичных профессиональных навыков преподавательской деятельности через проведение учебных занятий со студентами младших курсов.

Для достижения поставленной цели были определены следующие задачи:
\begin{itemize}
    \item ознакомиться со структурой и содержанием учебного процесса на кафедре;
    \item подготовить и провести цикл учебных занятий по дисциплине «Разработка кода информационных систем» для студентов СПО;
    \item освоить принципы планирования урока, подбора учебного материала и методов оценки знаний;
    \item приобрести практический опыт управления академической группой и поддержания учебной дисциплины.
\end{itemize}

Индивидуальное задание: в рамках учебной практики мне было поручено разработать и провести серию практических занятий по предмету «Разработка кода информационных систем» для студентов 3 курса СПО. Задание включало в себя создание практических лабораторных работ, направленных на освоение технологий HTML, CSS, JavaScript, работы с макетами Figma и построения UML-диаграмм, а также выполнение индивидуального домашнего задания, интегрирующего полученные навыки.

% ======================================================================
\section{Описание выполнения общей части программы практики}

В ходе прохождения учебной практики я был привлечен к проведению практических занятий по дисциплине «Разработка кода информационных систем» для студентов колледжа.

% ======================================================================
\section{Цель дисциплины}

Формирование у студентов системного представления о современных инструментальных средствах веб-разработки и выработка практических навыков их использования для создания статических и динамических веб-страниц, включая вёрстку макетов, добавление интерактивности и программную обработку данных.

Для достижения данной цели в рамках дисциплины решались следующие задачи:
\begin{itemize}
    \item Повышение учебной мотивации студентов через демонстрацию связи изучаемого материала с реальными задачами веб-разработки и карьерными возможностями в IT-сфере.
    \item Формирование понимания многообразия технологий веб-разработки: языков разметки, стилей, скриптовых языков и средств проектирования.
    \item Развитие способности к самостоятельному изучению и освоению новых инструментальных средств в зависимости от поставленной задачи.
    \item Выработка практических навыков работы с ключевыми инструментами современного веб-разработчика: редакторами кода, средами проектирования (Figma), браузерными инструментами разработчика, системами контроля версий.
\end{itemize}

% ======================================================================
\section{Содержание практических занятий}

Поскольку дисциплина со стороны магистра носила сугубо практико-ориентированный характер, учебный процесс был построен на выполнении цикла последовательных лабораторных работ, охватывающих путь веб-разработчика от основ вёрстки до создания динамических приложений с парсингом данных. Всего было разработано пять лабораторных работ и одно индивидуальное домашнее задание.

\begin{enumerate}
    \item \textbf{Лабораторная работа №1: «Основы языка разметки HTML. Вёрстка веб-страницы».} \\
    Задание: используя знания по HTML, выбрать любую страницу на Википедии и сверстать её. Страница должна быть воссоздана самостоятельно с помощью HTML-кода и визуально соответствовать оригинальной статье. Работа направлена на освоение структуры HTML-документа, основных тегов и их атрибутов.

    \item \textbf{Лабораторная работа №2: «Разработка статической веб-страницы. Вёрстка макета из Figma в HTML/CSS, создание диаграммы вариантов использования».} \\
    Задание: используя знания HTML и CSS, сверстать страницу по одному из макетов в Figma. Готовая веб-страница должна полностью соответствовать макету: соблюдена цветовая гамма, размеры и гарнитура шрифтов, отступы между блоками. Дополнительно построить диаграмму вариантов использования (Use-case Diagram).

    \item \textbf{Лабораторная работа №3: «Разработка статической веб-страницы. Вёрстка макета из Figma в HTML/CSS с применением стиля Grid, создание диаграмм кооперации и развертывания».} \\
    Задание: сверстать статическую веб-страницу по макету Figma с использованием технологии CSS Grid для организации структуры и расположения основных блоков. Дополнительно разработать UML-диаграмму кооперации и диаграмму развертывания.

    \item \textbf{Лабораторная работа №4: «Интерактивность и анимация на стороне клиента. Использование JavaScript для управления DOM, событиями и анимацией элементов веб-страницы».} \\
    Задание: взять готовую статическую страницу из лабораторной работы №2 и добавить сценарии на JavaScript так, чтобы страница демонстрировала осмысленное поведение: анимации запускаются по событиям, состояние элементов меняется предсказуемо. Реализовать минимум три различные по приёму анимации. Построить диаграмму состояний.

    \item \textbf{Лабораторная работа №5: «Разработка веб-страницы с парсингом данных на JavaScript».} \\
    Задание: разработать веб-страницу, которая осуществляет получение данных с публично доступного сайта или API (с использованием fetch API), и отображает эти данные в виде структурированного, визуально оформленного интерфейса. Дополнительно разработать диаграмму компонентов и диаграмму развертывания.

    \item \textbf{Индивидуальное домашнее задание: «Разработка веб-страницы с элементами анимации и обработкой данных. Вёрстка макета из Figma в HTML/CSS с интеграцией динамического контента».} \\
    Задание: сверстать веб-страницу по макету Figma, реализовать анимации с использованием CSS/JavaScript, а также механизм загрузки и обработки данных из локального JSON-файла с последующей подстановкой этих данных в элементы страницы. Построить три диаграммы: IDEF0, классовую диаграмму и диаграмму потоков данных (DFD).
\end{enumerate}

\subsection{Описание используемого методического материала и учебной литературы}

Для обеспечения высокого качества проведения занятий были разработаны технические задания к лабораторным работам, содержащие четкие пошаговые инструкции, теоретические справки и контрольные вопросы для самопроверки.

При подготовке материалов использовались как авторские наработки, так и актуальные общедоступные ресурсы, что позволило предоставить студентам наиболее релевантную и современную информацию:
\begin{itemize}
    \item \textbf{Wikipedia} — в качестве источника контента для вёрстки страницы в лабораторной работе №1.
    \item \textbf{Figma} — как основной инструмент для создания и экспорта графических макетов веб-страниц.
    \item \textbf{MDN Web Docs (Mozilla Developer Network)} — как справочный ресурс по HTML, CSS и JavaScript.
    \item \textbf{Документация по CSS Grid Layout и Flexbox} — для освоения современных методов позиционирования элементов.
    \item \textbf{Спецификация UML 2.5} — для изучения нотаций диаграмм (вариантов использования, кооперации, развертывания, состояний, компонентов, классов).
    \item \textbf{Стандарт IDEF0 и методология DFD} — для построения функциональных диаграмм и диаграмм потоков данных.
    \item \textbf{Публичные REST API (JSONPlaceholder, OpenWeatherMap, CoinGecko и др.)} — в качестве источников данных для парсинга.
    \item \textbf{Справочник по Web Animations API и requestAnimationFrame} — для реализации анимаций в лабораторной работе №4.
\end{itemize}

% ======================================================================
\section{Описание выполнения индивидуального задания}

В рамках учебной практики и в соответствии с индивидуальным заданием мною был разработан и апробирован комплекс из пяти последовательных лабораторных работ и одного индивидуального домашнего задания по дисциплине «Разработка кода информационных систем» для студентов колледжа. Данный комплекс представляет собой целостный учебно-методический продукт, направленный на системное погружение учащихся в современную экосистему веб-технологий и освоение ключевых инструментов веб-разработчика.

% ======================================================================
\section{Характеристика и цели разработанного комплекса}
\begin{sloppypar}
Основной целью разработки являлось создание практико-ориентированного курса, который бы не только знакомил студентов с разрозненными технологиями, а выстраивал у них целостное понимание процесса веб-разработки: от базовой вёрстки до создания интерактивных приложений с динамической обработкой данных. Комплекс лабораторных работ был спроектирован по принципу «от простого к сложному», что позволило последовательно решить следующие задачи:
\end{sloppypar}
\begin{enumerate}
    \item Сформировать фундаментальные навыки вёрстки веб-страниц с использованием HTML и CSS.
    \item Научить студентов работать с графическими макетами (Figma) и преобразовывать дизайн в код.
    \item Выработать практические навыки программирования на JavaScript для создания интерактивных веб-страниц.
    \item Освоить методы получения и обработки данных из внешних источников с помощью fetch API.
    \item Развить навыки системного проектирования через построение UML-диаграмм и функциональных моделей.
\end{enumerate}

\subsection{Подробный анализ разработанных лабораторных работ и их вклад в достижение целей практики}

\textbf{Лабораторная работа №1: «Основы языка разметки HTML. Вёрстка веб-страницы».}

\textit{Суть разработки:} разработано задание, в рамках которого студентам предлагается самостоятельно сверстать страницу из Википедии, воссоздав её структуру с использованием базовых тегов HTML: заголовков, абзацев, списков, таблиц, изображений и гиперссылок. Задание формирует понимание семантической структуры веб-документов.

\textit{Почему это важно для студента:} для учащегося эта работа является первой практической встречей с веб-разработкой. Она позволяет перевести теоретические знания о HTML в практическую плоскость, понять, как устроены веб-страницы «изнутри», и получить первый осязаемый результат — работающую веб-страницу.

\textit{Вклад в цели практики:} эта работа способствует формированию базовых технических навыков и повышению учебной мотивации, так как студент сразу видит результат своего труда и осознаёт, что способен создавать реальные веб-страницы.

\vspace{1em}
\textbf{Лабораторная работа №2: «Разработка статической веб-страницы. Вёрстка макета из Figma в HTML/CSS, создание диаграммы вариантов использования».}

\textit{Суть разработки:} разработано задание, целью которого является освоение процесса конвертации графического макета (дизайна) в работающий HTML/CSS-код. Студенты учатся работать с CSS-свойствами для позиционирования элементов, стилизации текста и создания визуально точных копий макетов. Дополнительно вводится понятие системного анализа через построение Use-case диаграммы.

\textit{Почему это важно для студента:} эта работа моделирует реальный рабочий процесс веб-разработчика: получение макета от дизайнера и его точная реализация в коде. Студент учится вниманию к деталям (цвета, шрифты, отступы) и приобретает навык вёрстки, востребованный в индустрии. Построение диаграммы вариантов использования развивает системное мышление и умение анализировать функциональные требования.

\textit{Вклад в цели практики:} это задание развивает профессиональные компетенции фронтенд-разработчика и закладывает основы проектного мышления. Как практикант, я отрабатывал навык конструирования учебного процесса, требующего от учащихся как технической точности, так и аналитической работы.

\vspace{1em}
\textbf{Лабораторная работа №3: «Вёрстка макета из Figma в HTML/CSS с применением стиля Grid, создание диаграмм кооперации и развертывания».}

\textit{Суть разработки:} эта работа углубляет навыки вёрстки через внедрение технологии CSS Grid — современного метода двумерного позиционирования элементов. Студенты осваивают грид-контейнеры, области и адаптивное поведение сетки. Дополнительно строятся диаграмма кооперации (взаимодействие компонентов при загрузке страницы) и диаграмма развертывания (физическое размещение компонентов).

\textit{Почему это важно для студента:} CSS Grid является индустриальным стандартом для построения сложных макетов. Освоение этой технологии даёт студенту мощный инструмент для создания современных веб-интерфейсов. Диаграммы кооперации и развертывания помогают понять архитектуру веб-приложений на системном уровне — как браузер, сервер и ресурсы взаимодействуют между собой.

\textit{Вклад в цели практики:} выполнение этого задания требовало от меня методического проектирования, структурирования обширного материала в доступную для усвоения форму, что является свидетельством развития навыков проектирования учебного процесса.

\vspace{1em}
\textbf{Лабораторная работа №4: «Интерактивность и анимация на стороне клиента. Использование JavaScript для управления DOM, событиями и анимацией».}

\textit{Суть разработки:} создано практическое задание, переводящее статическую страницу в интерактивную. Студенты осваивают базовый синтаксис JavaScript в контексте браузера: находят элементы на странице, реагируют на действия пользователя (клик, наведение, ввод), программно меняют классы, стили и содержимое. Реализуется минимум три анимации с использованием различных механизмов (CSS-классы, Web Animations API, requestAnimationFrame). Строится диаграмма состояний.

\textit{Почему это важно для студента:} это ключевой переход от статической вёрстки к динамической веб-разработке. Студент понимает, как «оживают» веб-страницы, как обрабатываются действия пользователя и как создаются визуальные эффекты. Диаграмма состояний учит формально описывать поведение системы — навык, необходимый при проектировании сложных интерфейсов.

\textit{Вклад в цели практики:} проводя эти занятия, я отрабатывал навыки организации учебного процесса и непосредственного взаимодействия с группой, объясняя концепции программирования и анимации на практике.

\vspace{1em}
\textbf{Лабораторная работа №5: «Разработка веб-страницы с парсингом данных на JavaScript».}

\textit{Суть разработки:} разработано задание на получение и обработку данных с внешних веб-ресурсов с помощью fetch API. Студенты выбирают публично доступный API (погода, курсы валют, каталоги фильмов и т.д.), реализуют асинхронную загрузку данных и их динамическое отображение на странице. Дополнительно строятся диаграмма компонентов (архитектура проекта) и диаграмма развертывания.

\textit{Почему это важно для студента:} эта работа знакомит с реальным паттерном современных веб-приложений: асинхронный запрос к серверу → обработка JSON → динамическое обновление DOM. Студент понимает, как работают сайты, отображающие актуальные данные (погода, новости, курсы валют). Диаграммы компонентов и развертывания закрепляют понимание архитектуры распределённых веб-систем.

\textit{Вклад в цели практики:} задание носит интегративный характер, объединяя навыки вёрстки, программирования и системного проектирования. Я, как практикант, отрабатывал навык составления комплексных заданий, проверяющих несколько компетенций одновременно.

\vspace{1em}
\textbf{Индивидуальное домашнее задание: «Вёрстка макета из Figma в HTML/CSS с интеграцией динамического контента».}

\textit{Суть разработки:} итоговое задание, интегрирующее все полученные навыки. Студент верстает страницу по макету Figma, добавляет анимации, реализует загрузку данных из локального JSON-файла и подстановку этих данных в элементы страницы. Строятся три диаграммы: IDEF0 (функциональная модель), классовая диаграмма и DFD (диаграмма потоков данных).

\textit{Почему это важно для студента:} это задание моделирует полноценный мини-проект, в котором студент выступает в роли full-stack веб-разработчика на клиентской стороне. Он проходит весь цикл: от дизайна до функционального прототипа с данными. Три диаграммы (IDEF0, Class, DFD) дают разносторонний взгляд на проектируемую систему — функциональный, структурный и потоковый.

\textit{Вклад в цели практики:} руководство выполнением этого комплексного задания позволило мне оценить эффективность всей методики преподавания и увидеть cumulative-эффект от последовательного прохождения лабораторных работ.

% ======================================================================
\section{Достижения студентов и результаты апробации}

В процессе проведения занятий отмечалась положительная динамика в освоении материала. Студенты, изначально не имевшие опыта веб-разработки, к завершению цикла:
\begin{itemize}
    \item Получили практические навыки вёрстки веб-страниц с использованием HTML и CSS.
    \item Освоили процесс преобразования графических макетов Figma в работающий код.
    \item Научились использовать CSS Grid и Flexbox для построения адаптивных макетов.
    \item Приобрели базовые навыки программирования на JavaScript: работа с DOM, обработка событий, создание анимаций.
    \item Получили опыт асинхронной загрузки и обработки данных из внешних API с помощью fetch.
    \item Научились строить и читать UML-диаграммы различных типов (Use-case, кооперации, развертывания, состояний, компонентов, классов), а также функциональные модели IDEF0 и DFD.
    \item Научились оформлять технические отчёты, структурировать информацию и объяснять технические концепции.
\end{itemize}

Особо следует отметить, что последовательная структура курса позволила студентам проследить эволюцию веб-страницы от статического HTML-документа до интерактивного приложения с динамической загрузкой данных, что способствовало формированию целостной картины процесса веб-разработки.

% ======================================================================
\section{Анализ педагогических трудностей и корректировка методики преподавания}

В ходе педагогической практики был выявлен ряд проблем, характерных для начального этапа обучения веб-технологиям, что потребовало гибкой корректировки методики проведения занятий:

\begin{enumerate}
    \item \textbf{Разный начальный уровень подготовки студентов.} Часть студентов имела базовое представление о HTML/CSS, в то время как другие сталкивались с этими технологиями впервые. Для решения этой проблемы был разработан дифференцированный подход: базовые инструкции дополнялись расширенными заданиями для более подготовленных студентов, а для начинающих проводились дополнительные мини-консультации.

    \item \textbf{Формальное использование цифровых инструментов.} Было замечено, что часть студентов использует генеративные нейросети для автоматического выполнения заданий, не вникая в содержание и не анализируя полученный код. В связи с этим была пересмотрена система оценивания: акцент смещён с количества предоставленного кода на его качественную защиту и устное пояснение. Введены задания, требующие не просто найти решение, а адаптировать его под конкретные требования макета.

    \item \textbf{Трудности с точностью вёрстки по макету.} Студенты часто пренебрегали точным соблюдением размеров, отступов и цветовой гаммы из Figma-макетов. Для решения этой проблемы были введены чек-листы проверки соответствия макету и практика взаимопроверки работ, что повысило как качество вёрстки, так и навыки code review.

    \item \textbf{Сложности при переходе от статики к динамике.} Переход от HTML/CSS к JavaScript (лабораторная работа №4) оказался наиболее трудным для студентов, не имевших опыта программирования. Для сглаживания этого перехода были добавлены дополнительные примеры кода с подробными комментариями и проведено занятие, посвящённое базовому синтаксису JavaScript.

    \item \textbf{Прокрастинация и низкая учебная мотивация.} Большинство студентов выполняло задания в «последний момент», что негативно сказывалось на глубине проработки материала. Для стимулирования ритмичной работы были введены промежуточные дедлайны и система бонусных баллов за своевременное выполнение.
\end{enumerate}

Таким образом, несмотря на выявленные трудности, апробация педагогических методик позволила адаптировать учебный процесс к реальным условиям и повысить вовлеченность студентов в итоговую проектную деятельность.

% ======================================================================
\section*{Заключение}
\addcontentsline{toc}{section}{Заключение}

В ходе практики мною был выполнен план проведения занятий и осуществлена полноценная разработка нового учебно-методического комплекса по дисциплине «Разработка кода информационных систем». Комплекс включает пять лабораторных работ, последовательно проводящих студента от основ HTML-вёрстки до создания интерактивных веб-приложений с парсингом данных, а также итоговое индивидуальное домашнее задание, интегрирующее все полученные навыки.

Данный комплекс носит практико-ориентированный характер, делая акцент на системном подходе к освоению веб-технологий: от статики к динамике, от вёрстки к программированию, от отдельных умений к комплексному проекту. Такой подход особенно актуален для студентов колледжа, стоящих в начале своего профессионального пути в IT-сфере.

Проделанная работа в полной мере соответствует целям и задачам учебной практики, способствуя развитию моих профессиональных навыков педагога и формируя прочную основу для будущей педагогической и профессиональной деятельности студентов.

% ======================================================================
\renewcommand{\refname}{Список использованных источников}
\addcontentsline{toc}{section}{Список использованных источников}

\begin{enumerate}
    \item MDN Web Docs: HTML, CSS, JavaScript справочник [Электронный ресурс] Mozilla.~--- Режим доступа: \url{https://developer.mozilla.org/ru/} (дата обращения: 05.06.2026).
    \item Wikipedia: Свободная энциклопедия [Электронный ресурс] ~--- Режим доступа: \url{https://wikipedia.org} (дата обращения: 05.06.2026).
    \item Figma: Collaborative Design Tool [Электронный ресурс] ~--- Режим доступа: \url{https://figma.com} (дата обращения: 05.06.2026).
    \item CSS Grid Layout: Спецификация [Электронный ресурс] W3C.~--- Режим доступа: \url{https://www.w3.org/TR/css-grid-1/} (дата обращения: 05.06.2026).
    \item CSS Flexible Box Layout: Спецификация [Электронный ресурс] W3C.~--- Режим доступа: \url{https://www.w3.org/TR/css-flexbox-1/} (дата обращения: 05.06.2026).
    \item UML 2.5 Specification [Электронный ресурс] Object Management Group.~--- Режим доступа: \url{https://www.omg.org/spec/UML/2.5} (дата обращения: 05.06.2026).
    \item Web Animations API [Электронный ресурс] W3C.~--- Режим доступа: \url{https://www.w3.org/TR/web-animations-1/} (дата обращения: 05.06.2026).
    \item Fetch API: Спецификация [Электронный ресурс] WHATWG.~--- Режим доступа: \url{https://fetch.spec.whatwg.org/} (дата обращения: 05.06.2026).
    \item JSONPlaceholder: Free Fake REST API [Электронный ресурс] ~--- Режим доступа: \url{https://jsonplaceholder.typicode.com} (дата обращения: 05.06.2026).
    \item IDEF0: Integration Definition for Function Modeling [Электронный ресурс] NIST.~--- Режим доступа: \url{https://www.idef.com/idef0} (дата обращения: 05.06.2026).
    \item Data Flow Diagram (DFD): Методология [Электронный ресурс] ~--- Режим доступа: \url{https://www.geeksforgeeks.org/what-is-dfd} (дата обращения: 05.06.2026).
    \item Roadmap.sh: Developer Roadmaps [Электронный ресурс] ~--- Режим доступа: \url{https://roadmap.sh} (дата обращения: 05.06.2026).
\end{enumerate}

% ======================================================================
\newpage
\section*{Приложение A. Лабораторная работа №1}
\addcontentsline{toc}{section}{Приложение A}

\textbf{Тема:} Основы языка разметки HTML. Вёрстка веб-страницы.

\textbf{Цель работы:} изучить основы языка разметки HTML, освоить структуру HTML-документа, познакомиться с основными тегами и их атрибутами. Получить практические навыки вёрстки веб-страницы на примере воссоздания статьи из Википедии.

\textbf{Задание:} используя знания по HTML, выбрать любую страницу на Википедии и сверстать её. Страница должна быть воссоздана самостоятельно с помощью HTML-кода и визуально соответствовать оригинальной статье.

Требования к выполнению:
\begin{itemize}
    \item Статья должна содержать заголовки разных уровней, абзацы, списки, таблицы, изображения, гиперссылки.
    \item HTML-код должен быть написан самостоятельно. Копирование исходного кода не допускается.
    \item Страница должна корректно отображаться в современных браузерах.
    \item Файл с расширением .html должен быть сохранён и предоставлен вместе с отчетом.
\end{itemize}

\section*{Приложение Б. Лабораторная работа №2}
\addcontentsline{toc}{section}{Приложение Б}

\textbf{Тема:} Разработка статической веб-страницы. Вёрстка макета из Figma в HTML/CSS, создание диаграммы вариантов использования.

\textbf{Цель работы:} изучить и применить на практике современные методы вёрстки веб-страниц. Освоить процесс конвертации графического макета в работающий HTML/CSS-код. Построить диаграмму вариантов использования.

\textbf{Задание:} сверстать страницу по макету Figma. Готовая страница должна полностью соответствовать макету по цветовой гамме, шрифтам, отступам. Построить Use-case диаграмму для своего варианта.

\section*{Приложение В. Лабораторная работа №3}
\addcontentsline{toc}{section}{Приложение В}

\textbf{Тема:} Вёрстка макета из Figma в HTML/CSS с применением CSS Grid, создание диаграмм кооперации и развертывания.

\textbf{Цель работы:} освоить преобразование макета Figma в код с использованием CSS Grid, приобрести навыки построения диаграмм кооперации и развертывания.

\textbf{Задание:} сверстать статическую страницу по макету Figma с применением CSS Grid. Разработать диаграмму кооперации (взаимодействие браузера, сервера и ресурсов) и диаграмму развертывания.

\section*{Приложение Г. Лабораторная работа №4}
\addcontentsline{toc}{section}{Приложение Г}

\textbf{Тема:} Интерактивность и анимация на стороне клиента. JavaScript для управления DOM, событиями и анимацией.

\textbf{Цель работы:} освоить базовый синтаксис JavaScript, научиться находить элементы в DOM, реагировать на события и программно менять классы, стили и содержимое. Реализовать минимум три анимации через различные механизмы.

\textbf{Задание:} взять страницу из ЛР №2 и добавить JavaScript-сценарии. Анимации запускаются по событиям, не менее двух разных типов событий. Построить диаграмму состояний.

\section*{Приложение Д. Лабораторная работа №5}
\addcontentsline{toc}{section}{Приложение Д}

\textbf{Тема:} Разработка веб-страницы с парсингом данных на JavaScript.

\textbf{Цель работы:} изучить методы получения и обработки данных с внешних ресурсов с помощью fetch API и JSON. Освоить динамическое обновление DOM.

\textbf{Задание:} разработать страницу, осуществляющую получение данных с публичного API и отображающую их в структурированном интерфейсе. Вёрстка с использованием CSS Grid или Flexbox. Построить диаграмму компонентов и диаграмму развертывания.

\section*{Приложение E. Индивидуальное домашнее задание}
\addcontentsline{toc}{section}{Приложение Е}

\textbf{Тема:} Разработка веб-страницы с элементами анимации и обработкой данных. Вёрстка макета из Figma в HTML/CSS с интеграцией динамического контента.

\textbf{Цель работы:} освоить преобразование макета Figma в код, приобрести навыки добавления анимаций и обработки данных из локальных источников с помощью JavaScript.

\textbf{Задание:} сверстать страницу по макету Figma, реализовать анимации (CSS/JS) и механизм загрузки данных из локального JSON-файла с подстановкой в элементы страницы. Построить диаграммы IDEF0, классовую диаграмму и DFD.

\end{document}
